\chapter{基础知识}

课件把本章分成五块：介绍，误差，浮点数与舍入误差，范数，线性方程组的性态与算法稳定性。
下面按同一顺序，每块先写对象与公式，再给一两个算例说明公式里的数代表什么。

\section{介绍}

用计算机解题的固定链条是
\[
\text{模型}\to\text{数值问题}\to\text{算法}\to\text{程序}.
\]
输出始终是近似值，因此后文只关心三件事：收敛、稳定、运算量。
本课遇到的数值问题可以写成三类：
\begin{align}
Ax&=b,\\
f(x)&=0,\\
\min_{\theta}\,L(\theta).
\end{align}
第一式输入 $A\in\mathbb{R}^{n\times n}$、$b\in\mathbb{R}^n$，输出 $x\in\mathbb{R}^n$，是未知量本身。
第二式输入非线性 $f$ 与初值，输出根 $x$。
第三式输入参数 $\theta$ 与损失 $L$，输出使 $L$ 变小的 $\theta$；$L(\theta)$ 是非负标量风险，不是预测值。

成立条件随问题而变：方形可逆组有唯一解；超定组 $m>n$ 一般无精确解，改成
\[
\min_x\,\lVert Ax-b\rVert_2^2
\quad\Leftrightarrow\quad
A^{\mathsf T}Ax=A^{\mathsf T}b;
\]
非线性与训练还要算法收敛。
形式上写 $x=A^{-1}b$ 并不要求真去求逆：求逆与乘法同为 $O(n^3)$，常数更大，实际用 LU、Cholesky、QR 或迭代直接得到一个 $x$。

\textbf{案例 A（线性组）。}
《九章算术》
\[
\begin{pmatrix}3&2&1\\2&3&1\\1&2&3\end{pmatrix}
\begin{pmatrix}x\\y\\z\end{pmatrix}
=\begin{pmatrix}39\\34\\26\end{pmatrix}
\]
的解是 $x=9.25$，$y=4.25$，$z=2.75$，单位为斗/秉。
这是 $Ax=b$ 的原型。

\textbf{案例 B（非线性）。}
Kepler 方程 $x-\varepsilon\sin x-t=0$（$0<\varepsilon<1$）对 $x$ 非线性，$x(t)$ 是偏近点角。
水温插值与人口拟合则分别是「穿过已知点」和「点多于参数时的最小二乘」：
人口用 $y\approx\beta_1 t^3+\beta_2 t^2+\beta_3 t+\beta_4$，输入 $\{(t_i,y_i)\}$，输出系数 $\beta$，再代入新 $t$ 得到预测人口。
神经网络把 $f_\theta$ 取得更宽，训练仍是 $\min_\theta L(\theta)$，梯度步
\[
\theta_{k+1}=\theta_k-\eta\nabla L(\theta_k)
\]
的 $\eta>0$ 是步长，$\nabla L$ 是损失对参数的梯度。

\textbf{案例 C（运算量）。}
朴素矩阵乘法 $(AB)_{ij}=\sum_{k=1}^n a_{ik}b_{kj}$ 要 $n^3$ 次乘法；
$n=10^5$、每秒 $10^9$ 次时约 $11.6$ 天。
Strassen 把指数从 $3$ 降到 $\log_2 7\approx 2.807$，说明同一问题可以换算法，但工程上仍受稳定性与缓存制约。

\section{误差}

数值分析假定模型与观测已给定，只跟踪两类计算误差：截断（方法）误差与舍入误差。
设 $x^*$ 为真值，$x$ 为近似值。绝对误差与相对误差为
\begin{equation}
\varepsilon=x-x^*,
\qquad
\xi=\frac{x-x^*}{x^*}.
\end{equation}
$\varepsilon$ 与 $x^*$ 同量纲，可正可负；$\lvert\varepsilon\rvert\le\varepsilon^*$ 时称 $\varepsilon^*$ 为绝对误差限，常写 $x^*=x\pm\varepsilon^*$。
$\xi$ 无量纲，$\lvert\xi\rvert=10^{-k}$ 大致对应 $k$ 位有效数字。
真值未知时改用 $\lvert x^*-x\rvert/\lvert x\rvert$。
$\lvert\varepsilon\rvert$ 小并不等于「够准」：同样差 $1$，对 $10^6$ 和对 $1$ 的意义不同，所以比较精度看 $\xi$。

有效数字：写 $x=\pm 0.a_1\ldots a_n\times 10^m$（$a_1\neq 0$），若
\[
\lvert x-x^*\rvert\le 0.5\times 10^{m-n},
\]
则 $x$ 有 $n$ 位有效数字。$n$ 是位数，不是 $x$ 的值。

误差沿函数传播。一元
\[
\Delta f(x^*)\approx\lvert f'(x^*)\rvert\,\lvert x-x^*\rvert,
\]
$\lvert f'(x^*)\rvert$ 是绝对误差的局部放大率。
二元一阶上界为
\[
\Delta f\approx\bigl\lvert\partial f/\partial u\bigr\rvert\Delta u
+\bigl\lvert\partial f/\partial v\bigr\rvert\Delta v.
\]
四则运算的估计误差：加减都是 $\Delta\tilde u+\Delta\tilde v$；
乘为 $\lvert\tilde u\rvert\Delta\tilde v+\lvert\tilde v\rvert\Delta\tilde u$；
除为同一分子除以 $\lvert\tilde v\rvert^2$。
减法的绝对误差限与加法相同，但真值接近时相对误差会被放大（相消）。

\bookfig{fig-p41-error-propagation.png}{第41页}{切线把自变量误差变成函数值误差的一阶估计。}{$\Delta\tilde x$：自变量绝对误差；估计误差：$\lvert f'\rvert\Delta\tilde x$；真误差：$\lvert f(x)-f(\tilde x)\rvert$}

\textbf{案例 A（有效数字）。}
$x^*=\pi$，$x=3.1415=0.31415\times 10^1$，故 $m=1$。
$\lvert\pi-3.1415\rvert\approx 9.3\times 10^{-5}$ 大于 $0.5\times 10^{1-5}$、小于 $0.5\times 10^{1-4}$，因此只有 $4$ 位有效数字。
四舍五入到四位：$42.195\to 42.20$，$0.0375551\to 0.03756$。

\textbf{案例 B（截断）。}
$I=\int_0^1 e^{-x^2}\,\mathrm{d}x$ 无初等原函数。逐项积分
\[
I=\sum_{k=0}^{\infty}\frac{(-1)^k}{k!\,(2k+1)},
\qquad
I_n=\sum_{k=0}^{n}\frac{(-1)^k}{k!\,(2k+1)},
\]
交错余项
\[
\lvert I-I_n\rvert\le\frac{1}{(n+1)!\,(2n+3)}.
\]
$I_n$ 是 $I$ 的近似（真值 $\approx 0.746824$）；右端是可事先计算的截断上界，不是随机误差。
$n=4$ 时 $I_4\approx 0.74749$，上界 $\approx 7.6\times 10^{-4}$，真误差 $\approx 6.6\times 10^{-4}$。

\section{浮点数及舍入误差}

浮点系统 $F(\beta,t,m,M)$ 表示
\[
x=W\beta^j,
\qquad
W=\pm\sum_{r=1}^{t}d_r\beta^{-r},
\quad -m\le j\le M.
\]
$\beta$ 是基，$t$ 是尾数位，$j$ 是阶码。规格化要求 $d_1\neq 0$（二进制即 $d_1=1$）。
$F$ 是有限散点集，不是 $\mathbb{R}$。
IEEE 单精度 $1+8+23$ 位，双精度 $1+11+52$ 位，数值为 $(-1)^S W 2^j$。

相邻规格化数间距 $\Delta=\beta^{j-t}$。四舍五入到最近浮点 $x_f$ 时
\[
\lvert x_f-x_R\rvert\le\tfrac12\beta^{j-t}.
\]
结合规格化下界 $\lvert x_R\rvert\ge\beta^{j-1}$，得到定理~1.3.1：
\begin{equation}
x_f=x_R(1+\delta),
\qquad
\lvert\delta\rvert\le u:=\tfrac12\beta^{1-t}.
\end{equation}
$\delta$ 是相对舍入，无量纲；$u$ 是单位舍入，只依赖 $\beta,t$，与 $\lvert x_R\rvert$ 无关（不溢出、规格化时）。
二进制 $u=2^{-t}$；双精度 $t=53$（含隐藏位）时 $u=2^{-53}\approx 1.11\times 10^{-16}$。

\bookfig{fig-p46-rounding-proof.png}{第46页}{由间距与规格化下界推出相对误差界 $u$。}{$\Delta$：相邻浮点间距；$u$：单位舍入；$\delta$：相对误差}

\section{向量范数与矩阵范数}

$\mathbb{R}^n$ 上的范数满足正定、绝对齐次、三角不等式。常用
\begin{align}
\lVert x\rVert_1&=\sum_i\lvert x_i\rvert,\\
\lVert x\rVert_2&=\bigl(\sum_i x_i^2\bigr)^{1/2},\\
\lVert x\rVert_\infty&=\max_i\lvert x_i\rvert.
\end{align}
输入向量，输出非负标量，表示长度。
有限维上一切范数等价：存在 $c_1,c_2>0$ 使 $c_1\lVert x\rVert'\le\lVert x\rVert\le c_2\lVert x\rVert'$，
因此「$x_k\to x^*$」与「$\lVert x_k-x^*\rVert\to 0$」对任意范数同时成立。
$\lVert x_k-x^*\rVert$ 是误差长度，不是某一个坐标。

矩阵范数另要求相容 $\lVert AB\rVert\le\lVert A\rVert\lVert B\rVert$。
Frobenius 范数把 $A$ 当成向量做 $2$-范数。
由向量范数诱导的算子范数
\[
\lVert A\rVert=\max_{x\neq 0}\frac{\lVert Ax\rVert}{\lVert x\rVert}
\]
是单位球上的最大拉伸，也是与该向量范数相容的最小矩阵范数，且 $\lVert I\rVert=1$。
列和、行和与谱范数分别为
\begin{align}
\lVert A\rVert_1&=\max_j\sum_i\lvert a_{ij}\rvert,\\
\lVert A\rVert_\infty&=\max_i\sum_j\lvert a_{ij}\rvert,\\
\lVert A\rVert_2&=\sqrt{\rho(A^{\mathsf T}A)}.
\end{align}
谱半径 $\rho(A)\le\lVert A\rVert$ 对任意矩阵范数成立；$\rho(A)$ 是按模最大特征值，控制 $A^k$ 是否趋于零。
若某算子范数下 $\lVert A\rVert<1$，则 $I\pm A$ 可逆，且
\[
\frac{1}{1+\lVert A\rVert}\le\lVert(I\pm A)^{-1}\rVert\le\frac{1}{1-\lVert A\rVert}.
\]

\textbf{案例。}
$x=(1,-2,3)^{\mathsf T}$：$\lVert x\rVert_1=6$，$\lVert x\rVert_2=\sqrt{14}$，$\lVert x\rVert_\infty=3$。
$A=\begin{pmatrix}1&-2\\-3&4\end{pmatrix}$：
$\lVert A\rVert_1=6$，$\lVert A\rVert_\infty=7$，
$\lVert A\rVert_2=\sqrt{15+\sqrt{221}}$。
对任意算子范数还有 $\lVert A^{-1}\rVert\ge 1/\lVert A\rVert$，以及
$\lVert A^{-1}-B^{-1}\rVert\le\lVert A^{-1}\rVert\lVert B^{-1}\rVert\lVert A-B\rVert$。

\section{线性方程组的性态，算法的稳定性}

解 $Ax=b$。右端扰动 $\delta b$ 给出
\begin{equation}
\frac{\lVert\delta x\rVert}{\lVert x\rVert}
\le
\operatorname{cond}(A)\,\frac{\lVert\delta b\rVert}{\lVert b\rVert},
\qquad
\operatorname{cond}(A)=\lVert A\rVert\lVert A^{-1}\rVert.
\end{equation}
$\operatorname{cond}(A)\ge 1$，无量纲，$\operatorname{cond}(kA)=\operatorname{cond}(A)$（$k\neq 0$）。
它是相对误差放大因子：输入相对右端扰动，输出相对解扰动的上界。
矩阵扰动 $\delta A$ 时还多一个分母 $1-\operatorname{cond}(A)\lVert\delta A\rVert/\lVert A\rVert$（需为正）。
$\operatorname{cond}(A)\gg 1$ 称病态，反之为良态。
经验上，行列式极小、行近似相关、主元过小、特征值数量级悬殊，都提示病态。
算法稳定性说的是舍入在计算过程中被放大多少；病态是问题固有的，稳定算法也不能凭空变出有效位。

\textbf{案例 A。}
\[
\begin{cases}x_1+x_2=2,\\ x_1+1.001 x_2=2,\end{cases}
\quad
\begin{cases}x_1+x_2=2,\\ x_1+1.001 x_2=2.001.\end{cases}
\]
解从 $(2,0)$ 变成 $(1,1)$。右端相对扰动约 $5\times 10^{-4}$，解的相对变化为 $O(1)$，放大约 $10^3$，与该 $2\times 2$ 矩阵的条件数同量级。

\textbf{案例 B（Hilbert）。}
\[
H_{ij}=\frac{1}{i+j-1},\quad i,j=1,\ldots,15,
\qquad
b=H\mathbf{1}.
\]
精确解 $x^*=\mathbf{1}$。双精度 $\operatorname{cond}_2(H)\approx 8.72\times 10^{17}$，
$u\cdot\operatorname{cond}_2(H)=O(10^2)$，有效数字可以全部丢失。
浮点解的残差 $\lVert Hx-b\rVert_2/\lVert b\rVert_2\approx 2\times 10^{-16}$，
相对误差 $\lVert x-x^*\rVert_2/\lVert x^*\rVert_2\approx 4.93$；
分量可偏离到 $-8.7$、$12.8$。
残差达到机器精度并不表示解对，这就是病态。
\begin{lstlisting}
import numpy as np
n, i = 15, np.arange(1, 16)
H = 1.0 / (i[:, None] + i[None, :] - 1.0)
x = np.linalg.solve(H, H @ np.ones(n))
\end{lstlisting}
\code{solve} 输出满足 $Hx\approx b$ 的 $x\in\mathbb{R}^{15}$，它不是 $x^*$。
